\chapter{Pare-feu pfSense}

\section{Pourquoi un pare-feu dans un lab AD ?}

Dans un environnement de lab isolé, un pare-feu n'est techniquement pas \textbf{obligatoire}. Les VMs communiquent sans problème sur le réseau NAT interne.

Cependant, \textbf{l'ajouter est crucial} pour deux raisons :

\begin{enumerate}
    \item \textbf{Réalisme professionnel} : Chaque déploiement AD en production possède un pare-feu frontal. Apprendre AD sans pare-feu revient à ignorer la moitié du métier.
    
    \item \textbf{Comprendre les ports AD} : Active Directory utilise une vingtaine de ports précis (53 DNS, 88 Kerberos, 389 LDAP, 445 SMB, etc.). Sans pare-feu, tout passe partout. Avec pfSense, si vous bloquez le port 88, Kerberos casse — vous apprenez \textit{pourquoi}.
\end{enumerate}

\section{Installation de pfSense}

\subsection{Téléchargement et création VM}

\begin{lstlisting}[language=bash]
# Télécharger pfSense CE (AMD64 ISO) depuis :
# https://www.netgate.com/pfsense-plus-software/try-pfsense

# Créer une VM dans VMware Workstation Pro :
# - Nom : fw-pfsense
# - Type : BSD, FreeBSD 64-bit
# - RAM : 1024 MB
# - Disque : 20 GB (VDI dynamique)
# - Carte 1 : NAT (accès Internet)
# - Carte 2 : Réseau interne AD-Lab
\end{lstlisting}

\subsection{Configuration console au démarrage}

Lors du premier démarrage, pfSense pose des questions :

\begin{lstlisting}[language=bash]
Should VLANs be set up now? [y|n] → n

Enter WAN interface name or 'a' for auto → em0
Enter LAN interface name or 'a' for auto → em1
Enter Optional interface 1 name → [Entrée — pas d'OPT1 pour maintenant]

Do you want to proceed? [y|n] → y

# À l'invite "Enter option :" tapez : 2
# Configure IP sur interface LAN :

Enter IPv4 address for LAN → 192.168.40.1
Enter IPv4 subnet bit count → 24
Enter upstream gateway (press Enter for none) → [Entrée]
Enable DHCP on LAN? → n
\end{lstlisting}

\section{Configuration initiale}

\subsection{Accès interface web}

Depuis une VM (ex : srv-dc01) ouvrez un navigateur :

\begin{lstlisting}
URL     : https://192.168.40.1
Login   : admin
Mot de passe : pfsense
\end{lstlisting}

Acceptez l'avertissement certificat SSL auto-signé.

\subsection{Assistant de configuration}

\begin{table}[H]
    \centering
    \begin{tabularx}{\textwidth}{|l|p{5cm}|}
        \hline
        \textbf{Paramètre} & \textbf{Valeur} \\
        \hline
        Hostname & fw-pfsense \\
        \hline
        Domain & corpmsi.local \\
        \hline
        Primary DNS & 192.168.40.10 \\
        \hline
        Secondary DNS & 192.168.40.11 \\
        \hline
        Timezone & Europe/Paris \\
        \hline
        NTP Server & 0.pfsense.pool.ntp.org \\
        \hline
        WAN Type & DHCP \\
        \hline
        LAN IP & 192.168.40.1/24 \\
        \hline
        Admin Password & DomainP@ssw0rd!Lab \\
        \hline
    \end{tabularx}
    \caption{Configuration pfSense}
    \label{tab:pfsense-config}
\end{table}

\section{Interfaces réseau}

\subsection{Assignation}

\begin{lstlisting}
Interfaces > Assignments
\end{lstlisting}

\begin{table}[H]
    \centering
    \begin{tabularx}{\textwidth}{|l|p{2cm}|p{4cm}|}
        \hline
        \textbf{Nom} & \textbf{Interface} & \textbf{Configuration} \\
        \hline
        WAN & em0 & DHCP (VMware NAT) \\
        \hline
        LAN & em1 & 192.168.40.1/24 \\
        \hline
    \end{tabularx}
    \caption{Interfaces pfSense}
    \label{tab:pfsense-interfaces}
\end{table}

\section{Règles de pare-feu}

\subsection{Principe général}

\begin{itemize}
    \item Les règles s'appliquent au trafic \textbf{entrant} sur chaque interface
    \item Première règle qui correspond s'applique (évaluation haut en bas)
    \item Trafic de retour autorisé automatiquement (stateful)
    \item Bloquer inter-VLAN utilisateurs
\end{itemize}

\subsection{Interface WAN}

\begin{lstlisting}[language=bash]
# Politique par défaut : BLOQUER TOUT

# Bloquer RFC1918 entrant (antispoofing)
Bloquer | RFC1918 (192.168.0.0/8, 172.16.0.0/12, 10.0.0.0/8)

# Bloquer bogons
Bloquer | Bogon networks

# Bloquer tout le reste (implicite)
\end{lstlisting}

\subsection{Interface LAN}

\begin{lstlisting}[language=bash]
# Laisser tout passer par défaut (lab)

# En production, ajouter des règles granulaires :
Autoriser | TCP/UDP → 53 (DNS)
Autoriser | TCP/UDP → 88 (Kerberos)
Autoriser | TCP → 389, 636 (LDAP)
Autoriser | TCP → 445 (SMB)
Autoriser | TCP → 3268, 3269 (GC)
Autoriser | TCP → 60000-60099 (RPC restreint)
Autoriser | TCP/UDP → 67, 68 (DHCP)
Autoriser | TCP/UDP → 123 (NTP)
\end{lstlisting}

\section{Ports Active Directory — Référence}

\begin{table}[H]
    \centering
    \begin{tabularx}{\textwidth}{|l|l|p{4cm}|}
        \hline
        \textbf{Port} & \textbf{Service} & \textbf{Si bloqué} \\
        \hline
        53 & DNS & Impossible joindre domaine \\
        \hline
        88 & Kerberos & "Cannot contact DC" \\
        \hline
        135 & RPC Mapper & GPOs n'appliquent pas \\
        \hline
        389 & LDAP & ADUC ne liste rien \\
        \hline
        445 & SMB & Scripts/profils ne chargent pas \\
        \hline
        3268 & Global Catalog & Recherche AD impossible \\
        \hline
        60000–60099 & RPC dynamique & Replication échoue \\
        \hline
    \end{tabularx}
    \caption{Ports AD critiques}
    \label{tab:ad-ports}
\end{table}

\section{DHCP Relay}

pfSense relaye les requêtes DHCP vers les DCs Windows (qui gèrent le DHCP).

\subsection{Configuration}

\begin{lstlisting}
Services > DHCP Relay
\end{lstlisting}

\begin{table}[H]
    \centering
    \begin{tabularx}{\textwidth}{|l|p{5cm}|}
        \hline
        \textbf{Paramètre} & \textbf{Valeur} \\
        \hline
        Enable DHCP Relay & Coché \\
        \hline
        Interfaces & LAN (sauf WAN) \\
        \hline
        DHCP Servers & 192.168.40.10 (srv-dc01), 192.168.40.11 (srv-dc02) \\
        \hline
    \end{tabularx}
    \caption{Configuration DHCP Relay}
    \label{tab:dhcp-relay}
\end{table}

\textbf{Important :} Désactiver le DHCP Server pfSense sur les interfaces. Le DHCP Relay et Server sont mutuellement exclusifs.

\section{NAT sortant}

\begin{lstlisting}
Firewall > NAT > Outbound
\end{lstlisting}

Mode : Automatic (convertit automatiquement le trafic LAN vers IP WAN)

\section{Troubleshooting pfSense}

\subsection{Pas de connectivité depuis les VMs}

\begin{lstlisting}[language=bash]
# Étape 1 : Vérifier la passerelle
ping 192.168.40.1

# Étape 2 : Vérifier accès WAN
ping 8.8.8.8

# Étape 3 : Vérifier DNS
nslookup microsoft.com 192.168.40.10

# Étape 4 : Vérifier si pfSense a Internet
# Depuis console pfSense : ping 8.8.8.8
\end{lstlisting}

\section{Checklist déploiement pfSense}

\begin{itemize}
    \item[\checkmark] Télécharger et installer pfSense
    \item[\checkmark] Configurer interfaces (em0 WAN, em1 LAN)
    \item[\checkmark] Configurer IP LAN (192.168.40.1/24)
    \item[\checkmark] Activer DHCP Relay vers srv-dc01 + srv-dc02
    \item[\checkmark] Désactiver DHCP Server pfSense
    \item[\checkmark] Configurer admin password
    \item[\checkmark] Ajouter forwarders DNS (8.8.8.8, 1.1.1.1)
    \item[\checkmark] Tester connectivité depuis vm
\end{itemize}
